% =====================================================================
% 09-conclusion.tex -- Conclusion and reproducibility.  E3, 3 min
% Sources: repository layout on main, scripts/20_train.py,
%          notebooks/kaggle_run.ipynb, results/run_log.csv
% =====================================================================
\renewcommand{\sectioninfo}{Presenter: E3 \quad $\cdot$ \quad 3 minutes}
\section{Conclusion and reproducibility}

% ------------------------------------------------------------------
\begin{frame}[fragile]{Reproduce everything from the repository\hfill\who{E3}}
  \begin{columns}[T]
    \begin{column}{0.47\textwidth}
\begin{lstlisting}[style=kdplain]
configs/    default.yaml, smoke.yaml
data/       teacher_sequences.jsonl
src/kd/     data, answer_extraction, losses,
            train, evaluate, stats
scripts/    10_cache_teacher_sequences.py
            20_train.py   30_evaluate.py
            40_make_figures.py
notebooks/  kaggle_run.ipynb, archive/
tests/      one pytest file per module
results/    run_log.csv, runs/<run_id>/
report/     note.qmd, references.bib, slides/
\end{lstlisting}
    \end{column}
    \begin{column}{0.5\textwidth}
\begin{lstlisting}[style=kdplain,language=bash,basicstyle=\ttfamily\tiny]
pip install -r requirements.txt && pip install -e .
CFG=configs/default.yaml
python scripts/10_cache_teacher_sequences.py --config $CFG
python scripts/20_train.py --config $CFG --arm onpolicy --seed 1
python scripts/30_evaluate.py --config $CFG --run onpolicy_s1
python -m kd.stats --out results/tables/results_table.csv
\end{lstlisting}
      \footnotesize
      \begin{itemize}
        \item \repofile{configs/smoke.yaml}: the same pipeline on 16 questions,
              about 15 minutes.
        \item Teacher answers and every run's predictions are committed, so the
              statistics rerun without a GPU.
      \end{itemize}
    \end{column}
  \end{columns}
  \begin{todobox}{README and figure script (E3)}
    \repofile{README.md} is still the 2-line stub and
    \repofile{scripts/40\_make\_figures.py} is empty on \repofile{main}; both are
    required by the brief.
  \end{todobox}
\end{frame}

% ------------------------------------------------------------------
\begin{frame}[fragile]{Every number traces back to code and settings\hfill\ours\ \who{E3}}
  \begin{columns}[T]
    \begin{column}{0.52\textwidth}
\begin{lstlisting}
# scripts/20_train.py records every run (shortened)
info = {
    "run_id": run_id, "arm": args.arm,
    "seed": args.seed, "started_utc": utc_now(),
    "git_commit": git_commit(),     # "+dirty" if edited
    "config_hash": file_hash(cfg),  # sha256 of the file
}
update_run_log(run_id, status="running", **info)
summary = run_training(cfg, arm, seed, out_dir)
# finally: status, end time, GPU-hours -> run_info.json
\end{lstlisting}
    \end{column}
    \begin{column}{0.45\textwidth}
      \footnotesize
      \begin{itemize}
        \item One row per run in \repofile{results/run\_log.csv}, one folder per run
              in \repofile{results/runs/}.
        \item All 8 main runs so far: commit \texttt{f2accc6}, config hash
              \texttt{85bee6cac023}: identical code and settings.
        \item On Kaggle, \repofile{notebooks/kaggle\_run.ipynb} only clones that
              commit and calls the scripts: no logic hidden in a notebook.
        \item Checkpoints stay on Kaggle; only small result files enter git.
      \end{itemize}
    \end{column}
  \end{columns}
\end{frame}

% ------------------------------------------------------------------
\begin{frame}{Conclusion\hfill\who{E3}}
  \begin{columns}[T]
    \begin{column}{0.55\textwidth}
      \begin{keybox}
        \small At equal updates and with the same target, \textbf{the student's own text}
        was the only source that made Qwen3-0.6B better at GSM8K: about
        \textbf{+7 points} over the untrained model in both seeds so far, while human
        and teacher text gave no gain. Matching the teacher's \textbf{probabilities}
        rather than hard labels was worth about \textbf{+16 points}.
      \end{keybox}
      \vspace{0.1cm}
      {\scriptsize\color{kdgray} Draft based on seeds 2 and 3; confirm with seed 1.}
    \end{column}
    \begin{column}{0.42\textwidth}
      \begin{block}{Next steps we did not run}
        \footnotesize
        \begin{itemize}
          \item A compute-matched comparison
          \item Reverse KL and mixed on/off-policy batches
          \item Other distillation temperatures
          \item A same-size teacher (Born-Again setting)
          \item A second dataset or model family
        \end{itemize}
      \end{block}
    \end{column}
  \end{columns}
  \vspace{0.15cm}
  \begin{todobox}{Final take-home message (E3 with the team)}
    Agree on one sentence the audience should remember, after seed 1 and the interpretation meeting.
  \end{todobox}
\end{frame}

% ------------------------------------------------------------------
\begin{frame}[plain,noframenumbering]
  \begin{tikzpicture}[remember picture,overlay]
    \fill[kdpaler] (current page.north west) rectangle (current page.south east);
    \fill[kdblue] (current page.north west) rectangle ([xshift=0.22cm]current page.south west);
  \end{tikzpicture}%
  \centering
  \vspace{1.2cm}
  {\color{kdnavy}\Huge\bfseries Questions?\par}
  \vspace{0.5cm}
  {\color{kdblue}\large Whose text should a small model practise on?\par}
  \vspace{0.8cm}
  {\small\color{kdgray} Code, results and this deck:
  \texttt{github.com/redkitex/Knowledge-Distillation-and-Model-Compression}\par}
\end{frame}
